% Options for packages loaded elsewhere
\PassOptionsToPackage{unicode}{hyperref}
\PassOptionsToPackage{hyphens}{url}
\PassOptionsToPackage{dvipsnames,svgnames,x11names}{xcolor}
\documentclass[
  spanish,
  12pt,
]{article}
\usepackage{xcolor}
\usepackage[letterpaper,left=2.25cm,right=2.25cm,top=3cm,bottom=2.3cm,headheight=42pt,headsep=0.6cm]{geometry}
\usepackage{amsmath,amssymb}
\setcounter{secnumdepth}{-\maxdimen} % remove section numbering
\usepackage{iftex}
\ifPDFTeX
  \usepackage[T1]{fontenc}
  \usepackage[utf8]{inputenc}
  \usepackage{textcomp} % provide euro and other symbols
\else % if luatex or xetex
  \usepackage{unicode-math} % this also loads fontspec
  \defaultfontfeatures{Scale=MatchLowercase}
  \defaultfontfeatures[\rmfamily]{Ligatures=TeX,Scale=1}
\fi
\usepackage{lmodern}
\ifPDFTeX\else
  % xetex/luatex font selection
\fi
% Use upquote if available, for straight quotes in verbatim environments
\IfFileExists{upquote.sty}{\usepackage{upquote}}{}
\IfFileExists{microtype.sty}{% use microtype if available
  \usepackage[]{microtype}
  \UseMicrotypeSet[protrusion]{basicmath} % disable protrusion for tt fonts
}{}
\makeatletter
\@ifundefined{KOMAClassName}{% if non-KOMA class
  \IfFileExists{parskip.sty}{%
    \usepackage{parskip}
  }{% else
    \setlength{\parindent}{0pt}
    \setlength{\parskip}{6pt plus 2pt minus 1pt}}
}{% if KOMA class
  \KOMAoptions{parskip=half}}
\makeatother
\usepackage{color}
\usepackage{fancyvrb}
\newcommand{\VerbBar}{|}
\newcommand{\VERB}{\Verb[commandchars=\\\{\}]}
\DefineVerbatimEnvironment{Highlighting}{Verbatim}{commandchars=\\\{\},fontsize=\footnotesize,breaklines=true}
% Add ',fontsize=\small' for more characters per line
\newenvironment{Shaded}{}{}
\newcommand{\AlertTok}[1]{\textcolor[rgb]{1.00,0.00,0.00}{\textbf{#1}}}
\newcommand{\AnnotationTok}[1]{\textcolor[rgb]{0.38,0.63,0.69}{\textbf{\textit{#1}}}}
\newcommand{\AttributeTok}[1]{\textcolor[rgb]{0.49,0.56,0.16}{#1}}
\newcommand{\BaseNTok}[1]{\textcolor[rgb]{0.25,0.63,0.44}{#1}}
\newcommand{\BuiltInTok}[1]{\textcolor[rgb]{0.00,0.50,0.00}{#1}}
\newcommand{\CharTok}[1]{\textcolor[rgb]{0.25,0.44,0.63}{#1}}
\newcommand{\CommentTok}[1]{\textcolor[rgb]{0.38,0.63,0.69}{\textit{#1}}}
\newcommand{\CommentVarTok}[1]{\textcolor[rgb]{0.38,0.63,0.69}{\textbf{\textit{#1}}}}
\newcommand{\ConstantTok}[1]{\textcolor[rgb]{0.53,0.00,0.00}{#1}}
\newcommand{\ControlFlowTok}[1]{\textcolor[rgb]{0.00,0.44,0.13}{\textbf{#1}}}
\newcommand{\DataTypeTok}[1]{\textcolor[rgb]{0.56,0.13,0.00}{#1}}
\newcommand{\DecValTok}[1]{\textcolor[rgb]{0.25,0.63,0.44}{#1}}
\newcommand{\DocumentationTok}[1]{\textcolor[rgb]{0.73,0.13,0.13}{\textit{#1}}}
\newcommand{\ErrorTok}[1]{\textcolor[rgb]{1.00,0.00,0.00}{\textbf{#1}}}
\newcommand{\ExtensionTok}[1]{#1}
\newcommand{\FloatTok}[1]{\textcolor[rgb]{0.25,0.63,0.44}{#1}}
\newcommand{\FunctionTok}[1]{\textcolor[rgb]{0.02,0.16,0.49}{#1}}
\newcommand{\ImportTok}[1]{\textcolor[rgb]{0.00,0.50,0.00}{\textbf{#1}}}
\newcommand{\InformationTok}[1]{\textcolor[rgb]{0.38,0.63,0.69}{\textbf{\textit{#1}}}}
\newcommand{\KeywordTok}[1]{\textcolor[rgb]{0.00,0.44,0.13}{\textbf{#1}}}
\newcommand{\NormalTok}[1]{#1}
\newcommand{\OperatorTok}[1]{\textcolor[rgb]{0.40,0.40,0.40}{#1}}
\newcommand{\OtherTok}[1]{\textcolor[rgb]{0.00,0.44,0.13}{#1}}
\newcommand{\PreprocessorTok}[1]{\textcolor[rgb]{0.74,0.48,0.00}{#1}}
\newcommand{\RegionMarkerTok}[1]{#1}
\newcommand{\SpecialCharTok}[1]{\textcolor[rgb]{0.25,0.44,0.63}{#1}}
\newcommand{\SpecialStringTok}[1]{\textcolor[rgb]{0.73,0.40,0.53}{#1}}
\newcommand{\StringTok}[1]{\textcolor[rgb]{0.25,0.44,0.63}{#1}}
\newcommand{\VariableTok}[1]{\textcolor[rgb]{0.10,0.09,0.49}{#1}}
\newcommand{\VerbatimStringTok}[1]{\textcolor[rgb]{0.25,0.44,0.63}{#1}}
\newcommand{\WarningTok}[1]{\textcolor[rgb]{0.38,0.63,0.69}{\textbf{\textit{#1}}}}
\usepackage{longtable,booktabs,array}
\newcounter{none} % for unnumbered tables
\usepackage{calc} % for calculating minipage widths
% Correct order of tables after \paragraph or \subparagraph
\usepackage{etoolbox}
\makeatletter
\patchcmd\longtable{\par}{\if@noskipsec\mbox{}\fi\par}{}{}
\makeatother
% Allow footnotes in longtable head/foot
\IfFileExists{footnotehyper.sty}{\usepackage{footnotehyper}}{\usepackage{footnote}}
\makesavenoteenv{longtable}
\usepackage{graphicx}
\makeatletter
\newsavebox\pandoc@box
\newcommand*\pandocbounded[1]{% scales image to fit in text height/width
  \sbox\pandoc@box{#1}%
  \Gscale@div\@tempa{\textheight}{\dimexpr\ht\pandoc@box+\dp\pandoc@box\relax}%
  \Gscale@div\@tempb{\linewidth}{\wd\pandoc@box}%
  \ifdim\@tempb\p@<\@tempa\p@\let\@tempa\@tempb\fi% select the smaller of both
  \ifdim\@tempa\p@<\p@\scalebox{\@tempa}{\usebox\pandoc@box}%
  \else\usebox{\pandoc@box}%
  \fi%
}
% Set default figure placement to htbp
\def\fps@figure{htbp}
\makeatother
\ifLuaTeX
\usepackage[bidi=basic,shorthands=off]{babel}
\else
\usepackage[bidi=default,shorthands=off]{babel}
\fi
\ifLuaTeX
  \usepackage{selnolig} % disable illegal ligatures
\fi
\setlength{\emergencystretch}{3em} % prevent overfull lines
\providecommand{\tightlist}{%
  \setlength{\itemsep}{0pt}\setlength{\parskip}{0pt}}
% Símbolos Unicode (griegos, sub/superíndices, operadores) que Latin Modern no trae
\usepackage{newunicodechar}
\newunicodechar{α}{\ensuremath{\alpha}}
\newunicodechar{β}{\ensuremath{\beta}}
\newunicodechar{δ}{\ensuremath{\delta}}
\newunicodechar{η}{\ensuremath{\eta}}
\newunicodechar{π}{\ensuremath{\pi}}
\newunicodechar{ρ}{\ensuremath{\rho}}
\newunicodechar{τ}{\ensuremath{\tau}}
\newunicodechar{χ}{\ensuremath{\chi}}
\newunicodechar{∈}{\ensuremath{\in}}
\newunicodechar{∉}{\ensuremath{\notin}}
\newunicodechar{∝}{\ensuremath{\propto}}
\newunicodechar{≈}{\ensuremath{\approx}}
\newunicodechar{≠}{\ensuremath{\ne}}
\newunicodechar{≤}{\ensuremath{\le}}
\newunicodechar{≥}{\ensuremath{\ge}}
\newunicodechar{⊆}{\ensuremath{\subseteq}}
\newunicodechar{→}{\ensuremath{\to}}
\newunicodechar{∑}{\ensuremath{\sum}}
\newunicodechar{·}{\ensuremath{\cdot}}
\newunicodechar{−}{\ensuremath{-}}
\newunicodechar{∞}{\ensuremath{\infty}}
\newunicodechar{Δ}{\ensuremath{\Delta}}
\newunicodechar{×}{\ensuremath{\times}}
\newunicodechar{₀}{\ensuremath{_{0}}}
\newunicodechar{₁}{\ensuremath{_{1}}}
\newunicodechar{₂}{\ensuremath{_{2}}}
\newunicodechar{₊}{\ensuremath{_{+}}}
\newunicodechar{ₖ}{\ensuremath{_{k}}}
\newunicodechar{ₙ}{\ensuremath{_{n}}}
\newunicodechar{ᵢ}{\ensuremath{_{i}}}
\newunicodechar{⁰}{\ensuremath{^{0}}}
\newunicodechar{⁶}{\ensuremath{^{6}}}
\newunicodechar{⁷}{\ensuremath{^{7}}}
\newunicodechar{⁹}{\ensuremath{^{9}}}
\newunicodechar{⁻}{\ensuremath{^{-}}}
\newunicodechar{ⁿ}{\ensuremath{^{n}}}
\usepackage{bookmark}
\IfFileExists{xurl.sty}{\usepackage{xurl}}{} % add URL line breaks if available
\urlstyle{same}
\hypersetup{
  pdftitle={Informe técnico: Metaheurísticas para el Problema del Ciclo Hamiltoniano sobre el FHCP Challenge Set},
  pdflang={es},
  colorlinks=true,
  linkcolor={Maroon},
  filecolor={Maroon},
  citecolor={Blue},
  urlcolor={Blue},
  pdfcreator={LaTeX via pandoc}}

\title{Informe técnico: Metaheurísticas para el Problema del Ciclo
Hamiltoniano sobre el FHCP Challenge Set}
\author{}
\date{}


\usepackage{fancyhdr}
\usepackage{setspace}
\usepackage{float}
\setstretch{1.15}
\pagestyle{fancy}
\fancyhf{}
\fancyhead[L]{\includegraphics[height=1.05cm]{docs/UTYP.png}}
\fancyhead[C]{\scriptsize Fundamentos de Inteligencia Artifical\\Unidad 2: Algoritmos heurísticos}
\fancyhead[R]{\includegraphics[height=1.05cm]{docs/LogoUPV_2023.png}}
\fancyfoot[C]{\small\thepage}
\renewcommand{\headrulewidth}{0.4pt}
\AtBeginEnvironment{longtable}{\footnotesize\setlength{\tabcolsep}{3pt}\setstretch{1.0}}
\setlength{\emergencystretch}{4em}
\hypersetup{pdftitle={Unidad 2: Algoritmos heurísticos},pdfauthor={Equipo de Fundamentos de Inteligencia Artifical}}
\usepackage{fvextra}

\begin{document}
\begin{titlepage}
\thispagestyle{empty}
\centering
\begin{tabular*}{\textwidth}{@{\extracolsep{\fill}}lr}
\includegraphics[height=1.8cm]{docs/UTYP.png}&\includegraphics[height=1.8cm]{docs/LogoUPV_2023.png}
\end{tabular*}
\vspace{0.5cm}

{\Large\bfseries UNIVERSIDAD POLITÉCNICA DE VICTORIA\par}
\vspace{0.4cm}
{\large Fundamentos de Inteligencia Artifical\par}
\vspace{0.7cm}
{\LARGE\bfseries Algoritmos heurísticos para el problema del ciclo hamiltoniano\par}
\vspace{0.3cm}
{\large Práctica de aula invertida\par}
{\large Informe técnico y comparativo del FHCP Challenge Set\par}
\vspace{0.5cm}
{\large\bfseries Unidad 2\quad |\quad Cuatrimestre 7\par}
\vspace{0.6cm}
{\small QUE PRESENTAN\par}
\vspace{0.15cm}
\begin{tabular}{lc}
\toprule
\textbf{Alumno(a)}&\textbf{Matrícula}\\
\midrule
Jared de Jesús Olazarán López&2430022\\
Georgina Reta Limas&2430043\\
Dulce Yaretzi Martínez Pérez&2430038\\
Melissa Jazmín Torres Martínez&2430017\\
Hugo Guillermo De León Ruiz&2330419\\
\bottomrule
\end{tabular}
\vspace{0.6cm}

{\small MAESTRO A CARGO DE LA MATERIA\par}
{\large\bfseries Doctor Said Polanco Martagon\par}
\vfill
\begin{flushright}
Ciudad Victoria, Tamaulipas\\13 de noviembre de 2026
\end{flushright}
\end{titlepage}
\pagenumbering{roman}
\renewcommand{\contentsname}{Contenido temático}
\tableofcontents
\clearpage
\pagenumbering{arabic}

\section{Introducción y objetivos}
El problema del ciclo hamiltoniano exige construir una ruta cerrada que visite cada vértice una sola vez. Su dificultad no consiste únicamente en ordenar vértices: también debe existir cada arista que une dos posiciones consecutivas, incluida la última con la primera. Esta práctica estudia cómo distintas metaheurísticas buscan soluciones sin enumerar todas las permutaciones posibles.

El objetivo general es implementar y evaluar GA, SA, ACO, colonia de abejas, DE y PSO sobre instancias del FHCP Challenge Set. Los objetivos específicos son: fundamentar la representación y el costo; preservar permutaciones legales; validar independientemente los ciclos; comparar éxito, calidad, evaluaciones y tiempo; y explicar la brecha respecto a métodos especializados. El informe sigue las cuatro fases de la práctica original: investigación, definición del problema y dataset, implementación y portafolio de evidencias.

La teoría orienta las decisiones de diseño: una permutación evita repetir u omitir vértices; contar aristas faltantes permite evaluar rutas incompletas; y aceptar temporalmente peores candidatos, compartir información o renovar poblaciones permite explorar regiones distintas. Ninguna de estas estrategias garantiza resolver todas las instancias dentro de un presupuesto finito. Fallar en una corrida tampoco demuestra que un grafo no sea hamiltoniano.

Los resultados experimentales que siguen proceden de los archivos CSV e historiales ya incluidos en el repositorio. La revisión actual del programa se presenta por separado en la guía de ejecución; no se atribuyen nuevos tiempos a los experimentos históricos.
\section{Resumen}\label{resumen}

Se implementaron seis metaheurísticas (GA, SA, ACO, Colonia de Abejas,
DE y PSO), más una búsqueda aleatoria como línea base y tres variantes
mejoradas (SA con 2-opt, GA con recombinación de aristas y ACO con
búsqueda local 2-opt con y sin preprocesamiento). Todas cargan grafos
del FHCP Challenge Set en formato HCP y devuelven un arreglo de vértices
(el ciclo propuesto), que un validador independiente comprueba.

Se evaluaron sobre 4 instancias del FHCP (n = 66 a 1002) con 10 semillas
por algoritmo, un presupuesto de 100 000 evaluaciones de fitness y un
tope de 15 s por corrida. Se hizo además una ronda ``pesada'' de 60 s
por corrida sobre 6 instancias (hasta n = 3468).

\textbf{Resultado principal:} solo los algoritmos con búsqueda local
dirigida resolvieron una instancia del set. ACO + 2-opt encontró el
ciclo hamiltoniano de \texttt{graph1} en las 10/10 corridas (≈0.5 s,
≈530 hormigas) y SA con 2-opt lo encontró en 5/10 corridas con 60 s. En
\texttt{graph24} y \texttt{graph41}, ACO + 2-opt y SA (2-opt) llegan a
fitness 1: encuentran un \textbf{camino} hamiltoniano, pero sus extremos
no son adyacentes. DE y PSO con claves aleatorias apenas mejoran a la
búsqueda aleatoria. En comparación, el equipo ganador del reto oficial
(Cohen y Coudert) resolvió 985 de los 1001 grafos.

\begin{center}\rule{0.5\linewidth}{0.5pt}\end{center}

\section{1. Fundamentación teórica (Fase
1)}\label{fundamentaciuxf3n-teuxf3rica-fase-1}

\subsection{1.1 Grafos, caminos y ciclos
hamiltonianos}\label{grafos-caminos-y-ciclos-hamiltonianos}

Un \textbf{grafo no dirigido} es un par G = (V, E), donde V = \{v₁,
\ldots, vₙ\} es el conjunto de \textbf{vértices} y E ⊆ \{\{u, v\} : u, v
∈ V, u ≠ v\} el de \textbf{aristas}. El \textbf{grado} de un vértice es
su número de vecinos.

\begin{itemize}
\tightlist
\item
  Un \textbf{camino} es una sucesión de vértices distintos (v₁, v₂,
  \ldots, vₖ) en la que cada par consecutivo \{vᵢ, vᵢ₊₁\} ∈ E.
\item
  Un \textbf{ciclo} es un camino (v₁, \ldots, vₖ) con k ≥ 3 que además
  cumple \{vₖ, v₁\} ∈ E.
\item
  Un \textbf{camino hamiltoniano} es un camino que pasa por los n
  vértices exactamente una vez.
\item
  Un \textbf{ciclo hamiltoniano} es un ciclo que pasa por los n vértices
  exactamente una vez y regresa al origen. Por tanto tiene
  \textbf{longitud exacta n}: usa n vértices distintos y n aristas, n −
  1 entre consecutivos más la de cierre \{vₙ, v₁\}.
\end{itemize}

Formalmente, G es hamiltoniano si existe una permutación π de \{1,
\ldots, n\} tal que \{v\_π(i), v\_π(i+1)\} ∈ E para i = 1, \ldots, n − 1
y \{v\_π(n), v\_π(1)\} ∈ E.

\textbf{Complejidad.} Decidir si un grafo es hamiltoniano es
\textbf{NP-completo}: es uno de los 21 problemas de Karp (1972), y la
versión de búsqueda es NP-difícil. El número de ciclos candidatos es (n
− 1)!/2. Para el grafo más pequeño del set (n = 66) eso da del orden de
10⁹⁰ candidatos, así que la fuerza bruta es impracticable.

\textbf{Relación con el TSP.} El ciclo hamiltoniano es un caso
particular del Problema del Agente Viajero. Basta asignar costo 0 a las
aristas de G y costo 1 a los pares que no son aristas: G es hamiltoniano
si y solo si el tour óptimo cuesta 0. Esta reducción justifica la
función de fitness de la sección 2.2 y explica por qué los mejores
resolvedores del FHCP (Concorde, LKH) vienen del TSP.

\subsection{1.2 Representación del espacio de
búsqueda}\label{representaciuxf3n-del-espacio-de-buxfasqueda}

Elegir la representación es decisivo. Hay dos riesgos: generar
estructuras \textbf{ilegales} (vértices repetidos u omitidos) o
representaciones con mucha redundancia. Se consideraron estas opciones:

{\def\LTcaptype{none} % do not increment counter
\begin{longtable}[]{@{}
  >{\raggedright\arraybackslash}p{(\linewidth - 4\tabcolsep) * \real{0.3333}}
  >{\raggedright\arraybackslash}p{(\linewidth - 4\tabcolsep) * \real{0.3333}}
  >{\raggedright\arraybackslash}p{(\linewidth - 4\tabcolsep) * \real{0.3333}}@{}}
\toprule\noalign{}
\begin{minipage}[b]{\linewidth}\raggedright
Representación
\end{minipage} & \begin{minipage}[b]{\linewidth}\raggedright
Ventaja
\end{minipage} & \begin{minipage}[b]{\linewidth}\raggedright
Problema
\end{minipage} \\
\midrule\noalign{}
\endhead
\bottomrule\noalign{}
\endlastfoot
Lista de aristas / matriz binaria & Directa & La mayoría de los
individuos no son ciclos; requiere reparación \\
\textbf{Permutación de vértices} (elegida) & Siempre visita cada vértice
una vez & Puede usar ``aristas'' que no existen en G \\
Claves aleatorias (\emph{random keys}) & Permite operadores continuos
(DE, PSO) & Fuerte redundancia: muchos vectores dan la misma
permutación \\
\end{longtable}
}

La decisión fue representar \textbf{toda solución como una permutación
de 0..n−1}. Con ella, la restricción ``cada vértice exactamente una
vez'' se cumple por construcción. La única restricción que puede
violarse es ``cada par consecutivo es una arista'', y eso es justo lo
que mide el fitness. Para que nunca aparezcan cromosomas ilegales, todos
los operadores (swap, inversión 2-opt, cruza OX, recombinación de
aristas, decodificación SPV) son \textbf{cerrados sobre permutaciones}.
Las pruebas de \texttt{tests/} lo verifican.

\subsection{1.3 Función de evaluación y
heurísticas}\label{funciuxf3n-de-evaluaciuxf3n-y-heuruxedsticas}

Con la reducción al TSP, el fitness natural es el \textbf{número de
aristas faltantes} del ciclo cerrado:

f(π) = \textbar\{ i ∈ \{1..n\} : \{π(i), π(i mod n + 1)\} ∉ E
\}\textbar{}

\begin{itemize}
\tightlist
\item
  f(π) = 0 si y solo si π es un ciclo hamiltoniano (criterio de éxito).
\item
  f(π) = 1 significa que π contiene un \textbf{camino hamiltoniano}
  cuyos extremos no son adyacentes.
\item
  La \textbf{legalidad} se define como 1 − f/n: la fracción de aristas
  del ciclo propuesto que existen en G.
\end{itemize}

Se estudiaron estas heurísticas que guían la búsqueda hacia la meta:

\begin{itemize}
\tightlist
\item
  \textbf{Evaluación incremental (delta).} Un swap o una inversión 2-opt
  cambia como máximo 4 o 2 aristas, así que el cambio de fitness se
  calcula en O(1) en lugar de O(n) (\texttt{two\_opt\_delta},
  \texttt{swap\_delta}).
\item
  \textbf{Regla de Warnsdorff} (1823). Se prefiere el vecino con menos
  vecinos libres, porque es el que corre más riesgo de quedar aislado.
  ACO la usa como información heurística η.
\item
  \textbf{Rotaciones de Pósa} (1976). Dado un camino, una arista desde
  un extremo hacia un vértice interior permite invertir un tramo y
  cambiar el extremo sin perder aristas. La búsqueda local de ACO las
  usa como movimientos neutros (delta = 0) para salir de mesetas.
\item
  \textbf{Aristas forzadas.} Un vértice de grado 2 obliga a usar sus dos
  aristas, lo que permite eliminar otras aristas incompatibles. Es un
  preprocesamiento clásico en resolvedores del HCP
  (\texttt{src/core/preprocess.py}).
\end{itemize}

\subsection{1.4 Métricas de
evaluación}\label{muxe9tricas-de-evaluaciuxf3n}

\textbf{Métricas de desempeño algorítmico}

{\def\LTcaptype{none} % do not increment counter
\begin{longtable}[]{@{}
  >{\raggedright\arraybackslash}p{(\linewidth - 2\tabcolsep) * \real{0.5000}}
  >{\raggedright\arraybackslash}p{(\linewidth - 2\tabcolsep) * \real{0.5000}}@{}}
\toprule\noalign{}
\begin{minipage}[b]{\linewidth}\raggedright
Métrica
\end{minipage} & \begin{minipage}[b]{\linewidth}\raggedright
Definición en este proyecto
\end{minipage} \\
\midrule\noalign{}
\endhead
\bottomrule\noalign{}
\endlastfoot
Tasa de éxito & Corridas con f = 0 (validadas) / corridas totales. IC 95
\% de Wilson \\
Costo en iteraciones & 1 iteración = 1 evaluación de fitness, completa o
por delta. Así se comparan algoritmos con poblaciones de distinto
tamaño \\
Tiempo computacional & Segundos de reloj por corrida
(\texttt{time.perf\_counter}) \\
ERT & \emph{Expected Running Time}: evaluaciones totales / número de
éxitos (Auger y Hansen, 2005) \\
Calidad & Fitness final (mejor, mediana, media, desviación) y
legalidad \\
\end{longtable}
}

\textbf{Métricas estadísticas: precisión, recall y F1.} Para aplicarlas,
la búsqueda del ciclo se plantea como una tarea de clasificación en dos
niveles:

\begin{itemize}
\tightlist
\item
  \textbf{Nivel corrida.} Las 1001 instancias del FHCP son
  hamiltonianas, así que toda corrida es un caso positivo.

  \begin{itemize}
  \tightlist
  \item
    TP: el algoritmo reporta un ciclo y el validador independiente lo
    confirma.
  \item
    FP: reporta un ciclo que el validador rechaza.
  \item
    FN: no encuentra ciclo.
  \end{itemize}

  Entonces precisión = TP/(TP+FP), \textbf{recall = TP/(TP+FN) = tasa de
  éxito} y F1 = 2PR/(P+R).
\item
  \textbf{Nivel arista.} De las n aristas que propone el ciclo, son
  aciertos las que existen en G. La \textbf{precisión de aristas} es
  igual a la legalidad. El recall de aristas no se puede definir sin un
  ciclo de referencia único (una instancia puede tener muchos ciclos),
  así que no se reporta.
\end{itemize}

\texttt{experiments/analisis.py} calcula estas métricas y las escribe en
\path{metricas_clasificacion.csv}.

\textbf{Pruebas estadísticas.} Se usaron Kruskal-Wallis por instancia,
Mann-Whitney unilateral del mejor algoritmo contra cada uno de los demás
con corrección de Holm, y Friedman con rangos medios entre instancias.
Son pruebas no paramétricas, porque el fitness es discreto y está
acotado.

\subsection{1.5 El FHCP Challenge}\label{el-fhcp-challenge}

El \textbf{Flinders Hamiltonian Cycle Project} publicó el \emph{FHCP
Challenge Set}: \textbf{1001 grafos hamiltonianos} de 66 a 9528 vértices
(media ≈ 3000), diseñados para ser difíciles para las heurísticas
estándar. Los candidatos se probaron con Concorde, LKH y SLH, y se
incluyó cada grafo que resultó difícil para al menos dos de ellos
(Haythorpe, 2018). El reto estuvo abierto del 30/09/2015 al 30/09/2016,
con un premio de US\$1001.

\begin{itemize}
\tightlist
\item
  \textbf{Formato HCP} (entrada): encabezado \texttt{NAME},
  \texttt{COMMENT}, \texttt{TYPE\ :\ HCP}, \texttt{DIMENSION\ :\ n},
  \texttt{EDGE\_DATA\_FORMAT\ :\ EDGE\_LIST},
  \texttt{EDGE\_DATA\_SECTION}; después una arista ``u v'' por línea con
  vértices 1..n, y al final \texttt{-1} / \texttt{EOF}.
\item
  \textbf{Formato TOUR} (solución): \texttt{TYPE\ :\ TOUR},
  \texttt{DIMENSION}, \texttt{TOUR\_SECTION}; los n vértices en orden,
  \texttt{-1}, \texttt{EOF}. El proyecto escribe así los ciclos
  encontrados (\texttt{results/tours/}).
\item
  \textbf{Criterio de benchmarking del reto:} el \textbf{número de
  instancias para las que se entrega un ciclo hamiltoniano válido}. No
  se premiaban soluciones parciales ni el tiempo. Por eso, además de la
  tasa de éxito por corrida, el proyecto reporta ``instancias
  resueltas'' (al menos una corrida exitosa).
\end{itemize}

\begin{center}\rule{0.5\linewidth}{0.5pt}\end{center}

\section{2. Definición del problema y del modelo (Fases 2 y
4.2)}\label{definiciuxf3n-del-problema-y-del-modelo-fases-2-y-4.2}

\subsection{2.1 Definición algorítmica del grafo y los
vértices}\label{definiciuxf3n-algoruxedtmica-del-grafo-y-los-vuxe9rtices}

\begin{itemize}
\tightlist
\item
  \texttt{src/core/hcp\_io.py} lee los archivos \texttt{.hcp} y
  convierte las etiquetas 1..n del archivo a vértices internos 0..n−1.
  Ignora aristas duplicadas y lazos. También escribe archivos
  \texttt{.tour} en 1..n.
\item
  \texttt{src/core/graph.py}:
  \texttt{Instance(name,\ n,\ adj,\ m,\ sha256)}, donde
  \texttt{adj{[}v{]}} es un \texttt{frozenset} de vecinos. Así, la
  consulta ``¿\{u, v\} ∈ E?'' es O(1) en promedio, que es la operación
  más usada por todos los algoritmos.
\item
  \texttt{src/core/fitness.py}: \texttt{fitness()} cuenta las aristas
  faltantes y \texttt{validate()} es un \textbf{validador independiente}
  de los algoritmos. Verifica que el resultado sea una permutación (sin
  repetidos ni ausentes) y que tenga 0 aristas faltantes. Todo éxito
  reportado se re-valida con él.
\end{itemize}

\subsection{2.2 Representación y función de
fitness}\label{representaciuxf3n-y-funciuxf3n-de-fitness}

{\def\LTcaptype{none} % do not increment counter
\begin{longtable}[]{@{}
  >{\raggedright\arraybackslash}p{(\linewidth - 4\tabcolsep) * \real{0.3333}}
  >{\raggedright\arraybackslash}p{(\linewidth - 4\tabcolsep) * \real{0.3333}}
  >{\raggedright\arraybackslash}p{(\linewidth - 4\tabcolsep) * \real{0.3333}}@{}}
\toprule\noalign{}
\begin{minipage}[b]{\linewidth}\raggedright
Algoritmo
\end{minipage} & \begin{minipage}[b]{\linewidth}\raggedright
Representación
\end{minipage} & \begin{minipage}[b]{\linewidth}\raggedright
Cómo se obtiene la permutación
\end{minipage} \\
\midrule\noalign{}
\endhead
\bottomrule\noalign{}
\endlastfoot
GA, SA, Bee & Permutación (vértice 0 fijo en SA/Bee) & Directa \\
ACO & Ruta construida vértice a vértice & Directa: la hormiga recorre
aristas reales y salta solo si no tiene salida \\
DE, PSO & Vector real de n claves & Decodificación \textbf{SPV}
(\emph{Smallest Position Value}): ordenar los vértices por clave \\
\end{longtable}
}

En todos los casos, fitness = aristas faltantes del ciclo cerrado
(sección 1.3), que se \textbf{minimiza}.

\subsection{2.3 Interfaz común}\label{interfaz-comuxfan}

Todos los algoritmos exponen

\begin{verbatim}
run(instance, config, seed) -> Result(
    best_route, best_fitness, history, time, iterations)
\end{verbatim}

Un objeto \texttt{Budget} lleva la cuenta de evaluaciones y tiempo, y
aplica el mismo límite a todos. \texttt{experiments/resolver.py} es el
punto de entrada mínimo que pide la Fase 4.1: carga un grafo, ejecuta
una metaheurística, imprime el arreglo de vértices y guarda el
\texttt{.tour} si es válido.

\begin{center}\rule{0.5\linewidth}{0.5pt}\end{center}

\section{3. Implementación algorítmica (Fase
3)}\label{implementaciuxf3n-algoruxedtmica-fase-3}

Los valores entre corchetes son los parámetros finales, elegidos con
\texttt{experiments/tuning.py} (ver 4.3).

\subsection{\texorpdfstring{3.1 Algoritmo Genético (GA): base de control
(\texttt{src/algorithms/ga.py})}{3.1 Algoritmo Genético (GA): base de control (src/algorithms/ga.py)}}\label{algoritmo-genuxe9tico-ga-base-de-control-srcalgorithmsga.py}

Población de permutaciones, selección por torneo (k = 3), \textbf{cruza
OX} (\emph{Order Crossover}: copia un tramo de un padre y completa con
el orden relativo del otro), \textbf{mutación por inversión} de un tramo
y elitismo (2 individuos). {[}pop\_size = 50, pc = 0.9, pm = 0.3{]}

\textbf{Variante GA (ER):} cruza por \textbf{recombinación de aristas}
(Whitley et al., 1989), que hereda aristas en lugar de posiciones.

\subsection{\texorpdfstring{3.2 Recocido Simulado (SA)
(\texttt{src/algorithms/sa.py})}{3.2 Recocido Simulado (SA) (src/algorithms/sa.py)}}\label{recocido-simulado-sa-srcalgorithmssa.py}

Como pide la práctica, el vecino es un \textbf{intercambio aleatorio de
dos índices (swap)} del arreglo. La aceptación sigue el criterio de
Metropolis: un movimiento que empeora en δ se acepta con probabilidad P
= exp(−δ/T). El enfriamiento es geométrico y \textbf{lento}: T va de t₀
a t\_end a lo largo de todo el presupuesto, T =
t₀·(t\_end/t₀)\^{}progreso, y se recalcula cada 64 movimientos. El delta
se calcula en O(1). {[}swap: t₀ = 2.0, t\_end = 0.01{]}

Con t₀ = 2.0, al inicio se acepta el 61 \% de los movimientos que
empeoran en una arista (e\^{}(−1/2)) y el 14 \% de los que empeoran en
cuatro. La temperatura baja 200 veces a lo largo de las 100 000
evaluaciones, así que al final SA solo acepta mejoras o movimientos
neutros (ver 4.3).

\textbf{Variante SA (2-opt):} el vecino invierte un tramo, lo que cambia
solo 2 aristas en vez de 4.

{[}t₀ = 0.3, t\_end = 0.02{]}

\subsection{\texorpdfstring{3.3 Colonia de Hormigas (ACO)
(\texttt{src/algorithms/aco.py})}{3.3 Colonia de Hormigas (ACO) (src/algorithms/aco.py)}}\label{colonia-de-hormigas-aco-srcalgorithmsaco.py}

Es un esquema \textbf{MAX-MIN Ant System} (Stützle y Hoos, 2000). Cada
hormiga parte del vértice 0 y avanza \textbf{solo por aristas reales}
hacia vértices no visitados. Elige el siguiente vértice con probabilidad
∝ τ\textsuperscript{α·η}β, con η = 1/(1 + vecinos libres del candidato),
que es la regla de Warnsdorff. Si queda acorralada, salta a un vértice
no visitado al azar; esa arista falta y penaliza el fitness. La feromona
se evapora a razón ρ y se deposita en las aristas del mejor de la
iteración y del mejor global, dentro de los límites {[}τ\_min,
τ\_max{]}. Así se \textbf{refuerzan las aristas que conectan nodos
difíciles} en rutas prometedoras. {[}ants = 5, α = 1, β = 2, ρ = 0.1{]}

\textbf{Variantes:}

\begin{itemize}
\tightlist
\item
  \textbf{ACO + 2-opt (\texttt{aco\_ls}):} cada hormiga se refina con
  una búsqueda local 2-opt dirigida a sus aristas faltantes, más hasta
  100 movimientos neutros por rotaciones de Pósa.
\item
  \textbf{ACO + 2-opt + prepro. (\texttt{aco\_ls\_pre}):} además reduce
  el grafo con aristas forzadas (sección 1.3).
\end{itemize}

\subsection{\texorpdfstring{3.4 Colonia de Abejas (ABC discreta)
(\texttt{src/algorithms/bee.py})}{3.4 Colonia de Abejas (ABC discreta) (src/algorithms/bee.py)}}\label{colonia-de-abejas-abc-discreta-srcalgorithmsbee.py}

Las fuentes de alimento son permutaciones.

\begin{itemize}
\tightlist
\item
  \textbf{Empleadas:} aplican una variación (swap o inversión) a su
  fuente y la aceptan si no empeora.
\item
  \textbf{Observadoras:} eligen una fuente con probabilidad ∝ 1/(1 + f),
  es decir, siguen a las fuentes prometedoras.
\item
  \textbf{Exploradoras:} abandonan una fuente sin mejora estricta en
  \texttt{limit} intentos y la reemplazan por una permutación aleatoria.
\end{itemize}

{[}sources = 5, limit = 2000{]}

\subsection{\texorpdfstring{3.5 Evolución Diferencial (DE)
(\texttt{src/algorithms/de.py})}{3.5 Evolución Diferencial (DE) (src/algorithms/de.py)}}\label{evoluciuxf3n-diferencial-de-srcalgorithmsde.py}

La DE es continua por naturaleza. Para adaptarla a permutaciones se usan
\textbf{claves aleatorias} (Bean, 1994): cada individuo es un vector x ∈
{[}0,1{]}ⁿ y se decodifica con SPV. La mutación DE/rand/1 (v = x\_r1 +
F·(x\_r2 − x\_r3)) y la cruza binomial operan en el espacio continuo, y
la selección es voraz. Así nunca se generan permutaciones ilegales.
{[}pop\_size = 10, F = 0.7, CR = 0.9{]}

\subsection{\texorpdfstring{3.6 Enjambre de Partículas (PSO)
(\texttt{src/algorithms/pso.py})}{3.6 Enjambre de Partículas (PSO) (src/algorithms/pso.py)}}\label{enjambre-de-partuxedculas-pso-srcalgorithmspso.py}

Cada partícula es un vector de claves con velocidad, decodificado por
SPV (Tasgetiren et al., 2007). La velocidad se actualiza con la regla
estándar v ← w·v + c₁r₁(pbest − x) + c₂r₂(gbest − x), con topología
global y velocidad acotada. No hay líder: cada partícula solo sigue
reglas simples con su mejor posición y la del enjambre. {[}swarm = 50, w
= 0.7, c₁ = c₂ = 1.5, vmax = 0.5{]}

\subsection{3.7 Línea base}\label{luxednea-base}

\texttt{random\_search} muestrea permutaciones aleatorias. Sirve de
piso: un algoritmo que no la supera no aprende nada de la estructura del
grafo.

\begin{center}\rule{0.5\linewidth}{0.5pt}\end{center}

\section{4. Dataset y diseño
experimental}\label{dataset-y-diseuxf1o-experimental}

\subsection{4.1 Instancias}\label{instancias}

{\def\LTcaptype{none} % do not increment counter
\begin{longtable}[]{@{}
  >{\raggedright\arraybackslash}p{(\linewidth - 10\tabcolsep) * \real{0.1667}}
  >{\raggedright\arraybackslash}p{(\linewidth - 10\tabcolsep) * \real{0.1667}}
  >{\raggedright\arraybackslash}p{(\linewidth - 10\tabcolsep) * \real{0.1667}}
  >{\raggedright\arraybackslash}p{(\linewidth - 10\tabcolsep) * \real{0.1667}}
  >{\raggedright\arraybackslash}p{(\linewidth - 10\tabcolsep) * \real{0.1667}}
  >{\raggedright\arraybackslash}p{(\linewidth - 10\tabcolsep) * \real{0.1667}}@{}}
\toprule\noalign{}
\begin{minipage}[b]{\linewidth}\raggedright
Instancia
\end{minipage} & \begin{minipage}[b]{\linewidth}\raggedright
n
\end{minipage} & \begin{minipage}[b]{\linewidth}\raggedright
m
\end{minipage} & \begin{minipage}[b]{\linewidth}\raggedright
Grado mín./máx.
\end{minipage} & \begin{minipage}[b]{\linewidth}\raggedright
Vértices de grado 2
\end{minipage} & \begin{minipage}[b]{\linewidth}\raggedright
Aristas eliminadas por preproceso
\end{minipage} \\
\midrule\noalign{}
\endhead
\bottomrule\noalign{}
\endlastfoot
graph1 & 66 & 99 & 3 / 3 & 0 & 0 \\
graph24 & 198 & 297 & 3 / 3 & 0 & 0 \\
graph41 & 300 & 451 & 3 / 4 & 0 & 0 \\
graph172 & 1002 & 1575 & 2 / 9 & 340 & 10 (672 aristas forzadas) \\
graph300* & 1815 & 3177 & 2 / 154 & 606 & 10 (1211 aristas forzadas) \\
graph600* & 3468 & 5203 & 3 / 4 & 0 & 0 \\
\end{longtable}
}

* Solo en la ronda pesada.

\textbf{Procedencia y verificación.} Las instancias provienen del FHCP
Challenge Set: \texttt{data/FHCPCS.7z} trae los 1001 grafos y
\texttt{data/catalogo\_fhcp.csv} los cataloga con n, m, grados y sha256.
El registro local declara una copia pública del conjunto, mientras que la procedencia oficial del archivo comprimido no está certificada en la evidencia revisada. Los saltos de línea CRLF y LF pueden producir hashes distintos aun para el mismo grafo; una comparación reproducible debe normalizarlos y verificar las aristas. Por tanto, se mantiene pendiente documentar la descarga oficial y su comparación con la copia utilizada.

\subsection{4.2 Protocolo}\label{protocolo}

{\def\LTcaptype{none} % do not increment counter
\begin{longtable}[]{@{}
  >{\raggedright\arraybackslash}p{(\linewidth - 10\tabcolsep) * \real{0.1667}}
  >{\raggedright\arraybackslash}p{(\linewidth - 10\tabcolsep) * \real{0.1667}}
  >{\raggedright\arraybackslash}p{(\linewidth - 10\tabcolsep) * \real{0.1667}}
  >{\raggedright\arraybackslash}p{(\linewidth - 10\tabcolsep) * \real{0.1667}}
  >{\raggedright\arraybackslash}p{(\linewidth - 10\tabcolsep) * \real{0.1667}}
  >{\raggedright\arraybackslash}p{(\linewidth - 10\tabcolsep) * \real{0.1667}}@{}}
\toprule\noalign{}
\begin{minipage}[b]{\linewidth}\raggedright
Ronda
\end{minipage} & \begin{minipage}[b]{\linewidth}\raggedright
Instancias
\end{minipage} & \begin{minipage}[b]{\linewidth}\raggedright
Algoritmos
\end{minipage} & \begin{minipage}[b]{\linewidth}\raggedright
Semillas
\end{minipage} & \begin{minipage}[b]{\linewidth}\raggedright
Presupuesto
\end{minipage} & \begin{minipage}[b]{\linewidth}\raggedright
Archivo
\end{minipage} \\
\midrule\noalign{}
\endhead
\bottomrule\noalign{}
\endlastfoot
\textbf{Benchmark principal} & graph1, 24, 41, 172 & los 10 + aleatoria
& 0..9 & 100 000 evaluaciones, tope 15 s & \texttt{results/}\allowbreak\texttt{bench\_*} \\
Pesada & + graph300, 600 & ga, sa\_2opt, aco, bee, aco\_ls, aco\_ls\_pre
& 0..9 & 60 s (evaluaciones sin límite práctico) &
\texttt{results/}\allowbreak\texttt{pesada\_*} \\
Media & graph1, 24, 41, 172 & ga, sa\_2opt, aco\_ls & 0..4 & 30 s &
\texttt{results/}\allowbreak\texttt{media\_*} \\
\end{longtable}
}

\begin{itemize}
\tightlist
\item
  \textbf{Comparación justa:} 1 iteración = 1 evaluación de fitness en
  todos los algoritmos. El tope de tiempo impide que los algoritmos con
  evaluaciones caras (GA-ER, ACO, ACO+LS) excedan el presupuesto. En ACO
  + LS, una iteración incluye la búsqueda local, así que su límite
  efectivo es el tiempo.
\item
  \textbf{Validación:} todo tour con f = 0 se re-valida con
  \texttt{validate()}. La columna \texttt{validado} de la tabla resumen
  vale 1 en todas las filas: \textbf{0 falsos positivos}.
\item
  \textbf{Entornos:} el benchmark principal se corrió en Python 3.12.3
  (Linux x86\_64; las 40 corridas de SA (swap) se repitieron el
  12/11/2026 con los parámetros de 4.3, ver nota en 5.2) y las rondas
  pesada/media en Python 3.14.6 (Linux x86\_64, 8 núcleos, ejecución en
  paralelo). La ejecución paralela puede introducir competencia por recursos. Los tiempos requieren cautela incluso dentro de una ronda y no deben compararse directamente entre entornos distintos.
\end{itemize}

\subsection{4.3 Ajuste de parámetros}\label{ajuste-de-paruxe1metros}

Se hizo un barrido por rejilla (\texttt{experiments/tuning.py},
resultados en \texttt{results/tuning.csv}) sobre graph1 y graph24 con
las semillas 100--102, \textbf{distintas de las del benchmark} para no
sobreajustar. El presupuesto fue de 30 000 evaluaciones y 8 s. Métrica:
fitness final medio / n.

{\def\LTcaptype{none} % do not increment counter
\begin{longtable}[]{@{}lll@{}}
\toprule\noalign{}
Algoritmo & Mejores parámetros & Fitness normalizado \\
\midrule\noalign{}
\endhead
\bottomrule\noalign{}
\endlastfoot
ACO & ants = 5, β = 2.0 (empate) & 0.011 \\
SA (2-opt) & t₀ = 0.3, t\_end = 0.02 & 0.096 \\
GA & pop\_size = 50, pm = 0.3 & 0.243 \\
SA (swap) & t₀ = 2.0, t\_end = 0.01 (ver abajo) & 0.223 \\
Bee & sources = 5, limit = 2000 & 0.295 \\
DE & pop\_size = 10, F = 0.7, CR = 0.9 & 0.624 \\
PSO & swarm = 50, w = 0.7, vmax = 0.5 & 0.767 \\
\end{longtable}
}

\begin{itemize}
\tightlist
\item
  \textbf{ACO:} ants = 5, β = 2 empató con ants = 10, β = 5 (0.01094 los
  dos). Se eligió 5/2 porque con menos hormigas por iteración la
  feromona se actualiza más seguido con el mismo presupuesto, y β = 2 es
  el valor habitual en MMAS (Stützle y Hoos, 2000).
\item
  \textbf{SA (swap):} el primer barrido (30 000 evaluaciones) eligió t₀
  = 0.1, t\_end = 0.05. Con esa temperatura un empeoramiento de 1 se
  acepta con probabilidad e⁻¹⁰ ≈ 0.00005, así que SA se comportaba como
  un descenso que solo acepta mejoras, y no como el recocido que pide la
  práctica. Se repitió el ajuste con el presupuesto real del benchmark
  (100 000 evaluaciones, \texttt{results/tuning\_sa.csv}). Las
  temperaturas casi nulas siguen siendo ligeramente mejores (0.204
  contra 0.223), pero se eligió la mejor configuración con t₀ ≥ 1,
  \textbf{t₀ = 2.0 y t\_end = 0.01}, para que SA acepte de verdad
  soluciones peores al inicio y enfríe lentamente. La diferencia entre
  temperaturas es pequeña frente a la que produce el vecindario (sección
  7).
\end{itemize}

\begin{center}\rule{0.5\linewidth}{0.5pt}\end{center}

\section{5. Comparativa de resultados (núcleo de la
práctica)}\label{comparativa-de-resultados-nuxfacleo-de-la-pruxe1ctica}

\subsection{5.1 Tasa de éxito y calidad (benchmark
principal)}\label{tasa-de-uxe9xito-y-calidad-benchmark-principal}

Cada celda indica \textbf{éxitos/10 · fitness mediano · mejor fitness}
(aristas faltantes; 0 = ciclo hamiltoniano):

{\def\LTcaptype{none} % do not increment counter
\begin{longtable}[]{@{}
  >{\raggedright\arraybackslash}p{(\linewidth - 8\tabcolsep) * \real{0.2000}}
  >{\raggedright\arraybackslash}p{(\linewidth - 8\tabcolsep) * \real{0.2000}}
  >{\raggedright\arraybackslash}p{(\linewidth - 8\tabcolsep) * \real{0.2000}}
  >{\raggedright\arraybackslash}p{(\linewidth - 8\tabcolsep) * \real{0.2000}}
  >{\raggedright\arraybackslash}p{(\linewidth - 8\tabcolsep) * \real{0.2000}}@{}}
\toprule\noalign{}
\begin{minipage}[b]{\linewidth}\raggedright
Algoritmo
\end{minipage} & \begin{minipage}[b]{\linewidth}\raggedright
graph1 (n=66)
\end{minipage} & \begin{minipage}[b]{\linewidth}\raggedright
graph24 (n=198)
\end{minipage} & \begin{minipage}[b]{\linewidth}\raggedright
graph41 (n=300)
\end{minipage} & \begin{minipage}[b]{\linewidth}\raggedright
graph172 (n=1002)
\end{minipage} \\
\midrule\noalign{}
\endhead
\bottomrule\noalign{}
\endlastfoot
\textbf{ACO + 2-opt + prepro.} & \textbf{10/10 · 0 · 0} & 0/10 · 1 · 1 &
0/10 · 1 · 1 & 0/10 · 12 · 9 \\
\textbf{ACO + 2-opt} & \textbf{10/10 · 0 · 0} & 0/10 · 1 · 1 & 0/10 · 1
· 1 & 0/10 · 13.5 · 9 \\
ACO & 0/10 · 1 · 1 & 0/10 · 1.5 · 1 & 0/10 · 1 · 1 & 0/10 · 16 · 13 \\
GA (ER) & 0/10 · 4.5 · 4 & 0/10 · 47.5 · 43 & 0/10 · 154 · 145 & 0/10 ·
911.5 · 898 \\
GA & 0/10 · 5.5 · 3 & 0/10 · 31.5 · 29 & 0/10 · 81 · 77 & 0/10 · 604.5 ·
594 \\
SA (2-opt) & 0/10 · 1 · 1 & 0/10 · 16 · 11 & 0/10 · 42.5 · 37 & 0/10 ·
634 · 622 \\
SA (swap) & 0/10 · 10.5 · 7 & 0/10 · 60.5 · 57 & 0/10 · 114 · 110 & 0/10
· 680 · 671 \\
Bee Colony & 0/10 · 4 · 2 & 0/10 · 55.5 · 51 & 0/10 · 131 · 125 & 0/10 ·
844 · 832 \\
DE & 0/10 · 29 · 26 & 0/10 · 129.5 · 124 & 0/10 · 219 · 210 & 0/10 ·
900.5 · 888 \\
PSO & 0/10 · 44 · 38 & 0/10 · 173 · 164 & 0/10 · 274 · 270 & 0/10 · 971
· 964 \\
Aleatoria & 0/10 · 53.5 · 52 & 0/10 · 185.5 · 184 & 0/10 · 287 · 286 &
0/10 · 989 · 986 \\
\end{longtable}
}

El IC 95 \% de Wilson de la tasa de éxito es {[}0.72, 1.00{]} para 10/10
y {[}0.00, 0.28{]} para 0/10.

\begin{figure}[H]
\centering
\pandocbounded{\includegraphics[keepaspectratio,alt={Tasa de éxito}]{results/analisis/fig_tasa_exito.png}}
\caption{Tasa de éxito}
\end{figure}

\begin{figure}[H]
\centering
\pandocbounded{\includegraphics[keepaspectratio,alt={Distribución del fitness final}]{results/analisis/fig_fitness_boxplot.png}}
\caption{Distribución del fitness final}
\end{figure}

\subsection{5.2 Tiempo computacional}\label{tiempo-computacional}

Cada celda indica \textbf{tiempo medio por corrida / evaluaciones
medias}:

{\def\LTcaptype{none} % do not increment counter
\begin{longtable}[]{@{}
  >{\raggedright\arraybackslash}p{(\linewidth - 8\tabcolsep) * \real{0.2000}}
  >{\raggedright\arraybackslash}p{(\linewidth - 8\tabcolsep) * \real{0.2000}}
  >{\raggedright\arraybackslash}p{(\linewidth - 8\tabcolsep) * \real{0.2000}}
  >{\raggedright\arraybackslash}p{(\linewidth - 8\tabcolsep) * \real{0.2000}}
  >{\raggedright\arraybackslash}p{(\linewidth - 8\tabcolsep) * \real{0.2000}}@{}}
\toprule\noalign{}
\begin{minipage}[b]{\linewidth}\raggedright
Algoritmo
\end{minipage} & \begin{minipage}[b]{\linewidth}\raggedright
graph1
\end{minipage} & \begin{minipage}[b]{\linewidth}\raggedright
graph24
\end{minipage} & \begin{minipage}[b]{\linewidth}\raggedright
graph41
\end{minipage} & \begin{minipage}[b]{\linewidth}\raggedright
graph172
\end{minipage} \\
\midrule\noalign{}
\endhead
\bottomrule\noalign{}
\endlastfoot
ACO + 2-opt + prepro. & 0.48 s / 532 & 15.00 s / 7 103 & 15.00 s / 3 392
& 15.01 s / 554 \\
ACO + 2-opt & 0.49 s / 532 & 15.00 s / 7 331 & 15.00 s / 3 339 & 15.01 s
/ 544 \\
ACO & 6.75 s / 100 000 & 15.00 s / 72 750 & 15.00 s / 44 590 & 15.00 s /
11 460 \\
GA (ER) & 7.78 s / 100 000 & 15.00 s / 68 368 & 15.00 s / 40 432 & 15.00
s / 9 987 \\
GA & 1.16 s / 100 000 & 2.20 s / 100 000 & 3.52 s / 100 000 & 13.77 s /
100 000 \\
SA (2-opt) & 0.12 s / 100 000 & 0.11 s / 100 000 & 0.17 s / 100 000 &
0.43 s / 100 000 \\
SA (swap)† & 0.22 s / 100 000 & 0.21 s / 100 000 & 0.24 s / 100 000 &
0.25 s / 100 000 \\
Bee Colony & 0.14 s / 100 000 & 0.14 s / 100 000 & 0.20 s / 100 000 &
0.46 s / 100 000 \\
DE & 1.03 s / 100 000 & 1.96 s / 100 000 & 2.87 s / 100 000 & 11.72 s /
100 000 \\
PSO & 1.19 s / 100 000 & 2.17 s / 100 000 & 3.21 s / 100 000 & 12.99 s /
100 000 \\
Aleatoria & 1.19 s / 100 000 & 3.45 s / 100 000 & 5.84 s / 100 000 &
15.00 s / 61 505 \\
\end{longtable}
}

† Las corridas de SA (swap) se repitieron con los parámetros nuevos en
otra máquina (Python 3.12.3, Linux x86\_64, 12 núcleos, ejecución
secuencial; ver \texttt{results/bench\_entorno.txt}). Su tiempo no es
directamente comparable con el resto de la tabla; el número de
evaluaciones sí.

\begin{itemize}
\tightlist
\item
  \textbf{ERT} de ACO + 2-opt en graph1: \textbf{532.5 evaluaciones}
  (≈0.5 s). Es el único caso con éxito en el benchmark principal.
\item
  \textbf{Costo por evaluación.} SA y Bee usan delta en O(1) y hacen 100
  000 movimientos en 0.1--0.5 s. GA, DE y PSO evalúan permutaciones
  completas en O(n), unas 10--100 veces más caro. ACO construye cada
  ruta en O(n·grado) y con búsqueda local cada hormiga cuesta mucho más
  (\textasciitilde1 ms en graph1), pero cada evaluación vale mucho más.
\end{itemize}

\begin{figure}[H]
\centering
\pandocbounded{\includegraphics[keepaspectratio,alt={Tiempo}]{results/analisis/fig_tiempo.png}}
\caption{Tiempo}
\end{figure}

\subsection{5.3 Convergencia}\label{convergencia}

La figura muestra la mediana del mejor fitness frente a las evaluaciones
(escala log). ACO y ACO + 2-opt llegan casi a 0 con pocas evaluaciones,
porque construyen rutas sobre aristas reales. SA (2-opt) y Bee bajan
rápido y se estancan. SA (swap) se queda cerca de la búsqueda aleatoria
durante las primeras \textasciitilde10 000 evaluaciones, porque a
temperatura alta acepta muchos empeoramientos, y desciende cuando la
temperatura baja. DE y PSO apenas se separan de la búsqueda aleatoria.

\begin{figure}[H]
\centering
\pandocbounded{\includegraphics[keepaspectratio,alt={Convergencia}]{results/analisis/fig_convergencia.png}}
\caption{Convergencia}
\end{figure}

\subsection{5.4 Precisión, recall y F1}\label{precisiuxf3n-recall-y-f1}

Datos de \texttt{results/analisis/metricas\_clasificacion.csv}:

{\def\LTcaptype{none} % do not increment counter
\begin{longtable}[]{@{}
  >{\raggedright\arraybackslash}p{(\linewidth - 8\tabcolsep) * \real{0.2000}}
  >{\raggedright\arraybackslash}p{(\linewidth - 8\tabcolsep) * \real{0.2000}}
  >{\raggedright\arraybackslash}p{(\linewidth - 8\tabcolsep) * \real{0.2000}}
  >{\raggedright\arraybackslash}p{(\linewidth - 8\tabcolsep) * \real{0.2000}}
  >{\raggedright\arraybackslash}p{(\linewidth - 8\tabcolsep) * \real{0.2000}}@{}}
\toprule\noalign{}
\begin{minipage}[b]{\linewidth}\raggedright
Algoritmo
\end{minipage} & \begin{minipage}[b]{\linewidth}\raggedright
Precisión (corrida)
\end{minipage} & \begin{minipage}[b]{\linewidth}\raggedright
Recall = tasa de éxito
\end{minipage} & \begin{minipage}[b]{\linewidth}\raggedright
F1
\end{minipage} & \begin{minipage}[b]{\linewidth}\raggedright
Precisión de aristas: g1 / g24 / g41 / g172
\end{minipage} \\
\midrule\noalign{}
\endhead
\bottomrule\noalign{}
\endlastfoot
ACO + 2-opt + prepro. & 1.00 & g1: 1.00; resto: 0 & g1: 1.00; resto: 0 &
1.000 / 0.995 / 0.997 / 0.988 \\
ACO + 2-opt & 1.00 & g1: 1.00; resto: 0 & g1: 1.00; resto: 0 & 1.000 /
0.995 / 0.997 / 0.987 \\
ACO & n/d & 0 & 0 & 0.985 / 0.991 / 0.995 / 0.984 \\
SA (2-opt) & n/d & 0 & 0 & 0.983 / 0.921 / 0.859 / 0.366 \\
GA & n/d & 0 & 0 & 0.920 / 0.839 / 0.728 / 0.397 \\
Bee Colony & n/d & 0 & 0 & 0.938 / 0.722 / 0.565 / 0.159 \\
GA (ER) & n/d & 0 & 0 & 0.924 / 0.758 / 0.488 / 0.093 \\
SA (swap) & n/d & 0 & 0 & 0.852 / 0.694 / 0.618 / 0.320 \\
DE & n/d & 0 & 0 & 0.562 / 0.335 / 0.260 / 0.099 \\
PSO & n/d & 0 & 0 & 0.348 / 0.133 / 0.087 / 0.031 \\
Aleatoria & n/d & 0 & 0 & 0.191 / 0.064 / 0.043 / 0.013 \\
\end{longtable}
}

La precisión por corrida es 1.00 siempre que hay positivos reportados:
\textbf{no hubo ningún falso positivo}, porque todo ciclo reportado pasó
el validador independiente. Es ``n/d'' cuando el algoritmo nunca reportó
un ciclo. El recall, que coincide con la tasa de éxito, es la métrica
que separa a los algoritmos. La precisión de aristas mide qué tan cerca
se quedó cada algoritmo de un ciclo cuando falló.

\subsection{5.5 Significancia
estadística}\label{significancia-estaduxedstica}

\begin{itemize}
\tightlist
\item
  \textbf{Kruskal-Wallis} por instancia: p \textless{} 10⁻¹⁶ en las 4
  instancias (entre 4.7·10⁻¹⁷ y 9.0·10⁻¹⁷). Los algoritmos difieren.
\item
  \textbf{Mann-Whitney + Holm}, comparando el mejor (ACO + 2-opt +
  prepro.) con cada uno de los demás: la diferencia es significativa al
  5 \% frente a \textbf{todos} los algoritmos en las 4 instancias,
  excepto frente a ACO + 2-opt sin preprocesamiento (p\_Holm = 1.0, 1.0,
  1.0, 0.30). Ver \texttt{results/analisis/pruebas\_estadisticas.csv}.
\item
  \textbf{Friedman} entre instancias: χ² = 33.8, \textbf{p = 0.0001}.
  Rangos medios (1 = mejor):
\end{itemize}

{\def\LTcaptype{none} % do not increment counter
\begin{longtable}[]{@{}
  >{\raggedright\arraybackslash}p{(\linewidth - 18\tabcolsep) * \real{0.1000}}
  >{\raggedright\arraybackslash}p{(\linewidth - 18\tabcolsep) * \real{0.1000}}
  >{\raggedright\arraybackslash}p{(\linewidth - 18\tabcolsep) * \real{0.1000}}
  >{\raggedright\arraybackslash}p{(\linewidth - 18\tabcolsep) * \real{0.1000}}
  >{\raggedright\arraybackslash}p{(\linewidth - 18\tabcolsep) * \real{0.1000}}
  >{\raggedright\arraybackslash}p{(\linewidth - 18\tabcolsep) * \real{0.1000}}
  >{\raggedright\arraybackslash}p{(\linewidth - 18\tabcolsep) * \real{0.1000}}
  >{\raggedright\arraybackslash}p{(\linewidth - 18\tabcolsep) * \real{0.1000}}
  >{\raggedright\arraybackslash}p{(\linewidth - 18\tabcolsep) * \real{0.1000}}
  >{\raggedright\arraybackslash}p{(\linewidth - 18\tabcolsep) * \real{0.1000}}@{}}
\toprule\noalign{}
\begin{minipage}[b]{\linewidth}\raggedright
ACO +\allowbreak 2opt +\allowbreak pre
\end{minipage} & \begin{minipage}[b]{\linewidth}\raggedright
ACO +\allowbreak 2opt
\end{minipage} & \begin{minipage}[b]{\linewidth}\raggedright
ACO
\end{minipage} & \begin{minipage}[b]{\linewidth}\raggedright
SA (2-opt)
\end{minipage} & \begin{minipage}[b]{\linewidth}\raggedright
GA
\end{minipage} & \begin{minipage}[b]{\linewidth}\raggedright
Bee
\end{minipage} & \begin{minipage}[b]{\linewidth}\raggedright
SA (swap)
\end{minipage} & \begin{minipage}[b]{\linewidth}\raggedright
GA (ER)
\end{minipage} & \begin{minipage}[b]{\linewidth}\raggedright
DE
\end{minipage} & \begin{minipage}[b]{\linewidth}\raggedright
PSO
\end{minipage} \\
\midrule\noalign{}
\endhead
\bottomrule\noalign{}
\endlastfoot
1.38 & 1.62 & 3.00 & 4.25 & 5.25 & 6.50 & 7.00 & 7.25 & 8.75 & 10.00 \\
\end{longtable}
}

\subsection{5.6 Ronda pesada (60 s por
corrida)}\label{ronda-pesada-60-s-por-corrida}

Cada celda indica \textbf{éxitos/10 · fitness mediano · mejor fitness}:

{\def\LTcaptype{none} % do not increment counter
\begin{longtable}[]{@{}
  >{\raggedright\arraybackslash}p{(\linewidth - 12\tabcolsep) * \real{0.1428}}
  >{\raggedright\arraybackslash}p{(\linewidth - 12\tabcolsep) * \real{0.1428}}
  >{\raggedright\arraybackslash}p{(\linewidth - 12\tabcolsep) * \real{0.1428}}
  >{\raggedright\arraybackslash}p{(\linewidth - 12\tabcolsep) * \real{0.1428}}
  >{\raggedright\arraybackslash}p{(\linewidth - 12\tabcolsep) * \real{0.1428}}
  >{\raggedright\arraybackslash}p{(\linewidth - 12\tabcolsep) * \real{0.1428}}
  >{\raggedright\arraybackslash}p{(\linewidth - 12\tabcolsep) * \real{0.1428}}@{}}
\toprule\noalign{}
\begin{minipage}[b]{\linewidth}\raggedright
Algoritmo
\end{minipage} & \begin{minipage}[b]{\linewidth}\raggedright
graph1 (66)
\end{minipage} & \begin{minipage}[b]{\linewidth}\raggedright
graph24 (198)
\end{minipage} & \begin{minipage}[b]{\linewidth}\raggedright
graph41 (300)
\end{minipage} & \begin{minipage}[b]{\linewidth}\raggedright
graph172 (1002)
\end{minipage} & \begin{minipage}[b]{\linewidth}\raggedright
graph300 (1815)
\end{minipage} & \begin{minipage}[b]{\linewidth}\raggedright
graph600 (3468)
\end{minipage} \\
\midrule\noalign{}
\endhead
\bottomrule\noalign{}
\endlastfoot
ACO + 2-opt + prepro. & \textbf{10/10 · 0 · 0} & 0/10 · 1 · 1 & 0/10 · 1
· 1 & 0/10 · 7.5 · \textbf{4} & 0/10 · 70.5 · 65 & 0/10 · 170 · 166 \\
ACO + 2-opt & \textbf{10/10 · 0 · 0} & 0/10 · 1 · 1 & 0/10 · 1 · 1 &
0/10 · 6.5 · 5 & 0/10 · 69 · 58 & 0/10 · 170 · 166 \\
ACO & 0/10 · 1 · 1 & 0/10 · 1.5 · 1 & 0/10 · 1 · 1 & 0/10 · 16 · 13 &
0/10 · \textbf{33 · 18} & 0/10 · \textbf{145 · 134} \\
SA (2-opt) & \textbf{5/10 · 0.5 · 0} & 0/10 · 1 · 1 & 0/10 · 1 · 1 &
0/10 · 25.5 · 20 & 0/10 · 133.5 · 119 & 0/10 · 969 · 901 \\
GA & 0/10 · 4 · 3 & 0/10 · 18 · 15 & 0/10 · 28.5 · 24 & 0/10 · 563 · 551
& 0/10 · 1470.5 · 1421 & 0/10 · 3296.5 · 3286 \\
Bee Colony & 0/10 · 2 · 1 & 0/10 · 30 · 28 & 0/10 · 71.5 · 70 & 0/10 ·
549 · 526 & 0/10 · 1302 · 1262 & 0/10 · 3124 · 3028 \\
\end{longtable}
}

\begin{itemize}
\tightlist
\item
  Con 60 s, \textbf{SA (2-opt) también resuelve graph1} (5/10). Las 5
  corridas exitosas tardaron en promedio 39.7 s y 1.2·10⁷ movimientos;
  las otras 5 agotaron los 60 s. Es la segunda metaheurística del
  proyecto que encuentra un ciclo hamiltoniano del set.
\item
  En los grafos más grandes (graph300, graph600), \textbf{ACO sin
  búsqueda local supera a ACO + 2-opt}. Con n ≥ 1815, la búsqueda local
  deja solo \textasciitilde70--200 hormigas en 60 s y la feromona no
  alcanza a aprender. Friedman en esta ronda: χ² = 23.6, p = 0.0003.
  Análisis completo en \texttt{results/analisis\_pesada/}.
\end{itemize}

\begin{center}\rule{0.5\linewidth}{0.5pt}\end{center}

\section{6. Benchmarking contra el FHCP
Challenge}\label{benchmarking-contra-el-fhcp-challenge}

\subsection{6.1 Resultados oficiales del
reto}\label{resultados-oficiales-del-reto}

{\def\LTcaptype{none} % do not increment counter
\begin{longtable}[]{@{}ll@{}}
\toprule\noalign{}
Equipo & Grafos resueltos (de 1001) \\
\midrule\noalign{}
\endhead
\bottomrule\noalign{}
\endlastfoot
N. Cohen y D. Coudert (INRIA, Francia), ganadores & \textbf{985} \\
Johnson (IBM, Reino Unido) & 614 \\
Gharbi y Syarif (King Saud University) & 488 \\
Noisternig (TU Darmstadt) & 464 \\
Nurhafiz (independiente) & 385 \\
\end{longtable}
}

Fuente: Haythorpe (2018) y la página oficial del FHCP. Ningún equipo
resolvió los 1001. Los 16 grafos que quedaron sin resolver son
instancias derivadas de grafos de Stedman (Sted4, Sted5, Sted6 y 13
combinaciones que los contienen). Los ganadores descompusieron las
instancias, buscaron estructuras repetidas entre ellas y reutilizaron
soluciones de subgrafos. En la construcción del set se usaron como
referencia resolvedores especializados del TSP/HCP: Concorde (exacto,
\emph{branch-and-cut}) y LKH (heurística de Lin-Kernighan-Helsgaun).

\subsection{6.2 Contraste con este
proyecto}\label{contraste-con-este-proyecto}

{\def\LTcaptype{none} % do not increment counter
\begin{longtable}[]{@{}
  >{\raggedright\arraybackslash}p{(\linewidth - 4\tabcolsep) * \real{0.3333}}
  >{\raggedright\arraybackslash}p{(\linewidth - 4\tabcolsep) * \real{0.3333}}
  >{\raggedright\arraybackslash}p{(\linewidth - 4\tabcolsep) * \real{0.3333}}@{}}
\toprule\noalign{}
\begin{minipage}[b]{\linewidth}\raggedright
Aspecto
\end{minipage} & \begin{minipage}[b]{\linewidth}\raggedright
Mejores del FHCP
\end{minipage} & \begin{minipage}[b]{\linewidth}\raggedright
Este proyecto
\end{minipage} \\
\midrule\noalign{}
\endhead
\bottomrule\noalign{}
\endlastfoot
Instancias resueltas & 985/1001 (ganadores); 385--614 (resto de los
equipos) & \textbf{1} de 6 probadas (graph1, n = 66) \\
Tamaño máximo atacado & hasta 9528 vértices & 3468 vértices \\
Mejor resultado parcial & n/a (solo cuentan ciclos completos) & Camino
hamiltoniano (f = 1) en graph24 y graph41; f = 4 en graph172 \\
Técnicas & Descomposición estructural, reducción al TSP, Concorde/LKH,
búsqueda local profunda, días de cómputo & Metaheurísticas genéricas,
15--60 s por corrida, Python puro \\
\end{longtable}
}

\textbf{Análisis.}

\begin{enumerate}
\def\labelenumi{\arabic{enumi}.}
\tightlist
\item
  \textbf{Brecha de calidad.} Con la métrica oficial (instancias con
  ciclo válido), las metaheurísticas genéricas quedan muy lejos de los
  resolvedores especializados. Incluso graph1, la instancia más pequeña del set, solo se resolvió con búsqueda local dirigida (2-opt + rotaciones
  de Pósa). Es justo el tipo de componente que tienen LKH y los
  resolvedores ganadores.
\item
  \textbf{``Casi'' no cuenta.} En graph24 y graph41 los mejores
  algoritmos encuentran caminos hamiltonianos (99.5--99.7 \% de aristas
  correctas) pero no logran cerrar el ciclo. En estas corridas, la búsqueda alcanza rutas casi correctas, pero se estanca antes de cerrar el ciclo; el experimento no mide cuántos ciclos existen en cada grafo. Una heurística que
  solo mide aristas faltantes ve una meseta enorme en f = 1.
\item
  \textbf{Estructura.} Las instancias graph1 y graph24 son cúbicas (grado 3); graph41 incluye grado 4. Ninguna de esas tres tiene vértices de grado 2, así que el preprocesamiento de aristas
  forzadas no cambia nada en ellos. Por eso aco\_ls y aco\_ls\_pre dan
  resultados idénticos en graph1/24/41. En graph172 y graph300 fija 672
  y 1211 aristas, pero el efecto no es significativo (p\_Holm = 0.30 en
  graph172; en la ronda pesada el mejor f baja de 5 a 4, aunque la
  mediana sube de 6.5 a 7.5). Los ganadores explotaron esta clase de
  estructura a gran escala.
\item
  \textbf{Escala.} El costo de la búsqueda local en Python limita a ACO
  + 2-opt en n ≥ 1000. Una implementación compilada podría reducir este costo; su mejora concreta debe medirse experimentalmente.
\end{enumerate}

\begin{center}\rule{0.5\linewidth}{0.5pt}\end{center}

\section{7. Conclusiones}\label{conclusiones}

\begin{enumerate}
\def\labelenumi{\arabic{enumi}.}
\tightlist
\item
  \textbf{ACO fue la metaheurística que mejor se adaptó a los cuellos de
  botella de los grafos difíciles.} Construye las rutas \textbf{sobre
  aristas reales} con la regla de Warnsdorff, así que nunca propone
  aristas inexistentes salvo en callejones sin salida, y la feromona
  aprende qué aristas conectan los nodos difíciles. Sola, llega a f = 1
  en graph1/24/41. Con búsqueda local 2-opt + Pósa, \textbf{resuelve
  graph1 en 10/10 corridas en medio segundo} y queda primera en el
  ranking de Friedman (p = 0.0001).
\item
  \textbf{SA es la segunda mejor opción si se usa el vecindario
  adecuado.} SA con swap, como lo pide el enunciado, cambia 4 aristas
  por movimiento y se estanca. SA con 2-opt cambia solo 2: reduce el
  fitness mediano entre 2.7 y 10.5 veces en n ≤ 300 (7 \% en graph172) y
  con 60 s resolvió graph1 en 5/10 corridas. La representación y el
  vecindario pesan más que el esquema de enfriamiento: cambiar la
  temperatura de SA (swap) movió el fitness ≈10 \% (sección 4.3),
  mientras que cambiar el vecindario lo movió hasta 10 veces.
\item
  \textbf{GA y Bee Colony} mejoran claramente a la búsqueda aleatoria,
  pero se degradan rápido con n.~La cruza OX conserva el orden pero
  rompe aristas. La recombinación de aristas (GA-ER) es mejor por
  evaluación, pero cuesta 7 veces más y con tope de tiempo quedó peor.
\item
  \textbf{DE y PSO con claves aleatorias no son adecuados para este
  problema.} El mapeo continuo → permutación (SPV) destruye la
  localidad: un cambio pequeño en las claves reordena muchos vértices.
  Por eso su fitness queda cerca de la búsqueda aleatoria, sobre todo en
  n = 1002.
\item
  \textbf{Lección central.} Ante grafos ``engañosos'', lo que funcionó
  fue \textbf{explotar la estructura del grafo}: construir solo sobre
  aristas reales, búsqueda local dirigida a las aristas faltantes,
  movimientos neutros para cruzar mesetas y aristas forzadas. Ajustar
  más los parámetros de un algoritmo ciego a la estructura no habría
  dado ese resultado. Es la misma conclusión que se desprende de los
  resultados del reto oficial.
\end{enumerate}

\textbf{Trabajo futuro:} usar un fitness que distinga caminos según qué
tan cerca estén sus extremos de cerrarse; añadir movimientos 3-opt /
Or-opt; reiniciar la búsqueda desde varios extremos con rotaciones de
Pósa; reimplementar el núcleo en C/Numba; y probar sobre un subconjunto
mayor del set (el catálogo ya permite seleccionar por n).

\begin{center}\rule{0.5\linewidth}{0.5pt}\end{center}

\section{8. Reproducibilidad}\label{reproducibilidad}

\begin{Shaded}
\begin{Highlighting}[]
\ExtensionTok{pip}\NormalTok{ install }\AttributeTok{{-}r}\NormalTok{ requirements.txt}
\ExtensionTok{python} \AttributeTok{{-}m}\NormalTok{ pytest                          }\CommentTok{\# 48 pruebas pasan, 1 omitida}
\CommentTok{\# Extraer data/FHCPCS.7z en data/instances/}

\CommentTok{\# Resolver una instancia (devuelve el arreglo de vértices y guarda el .tour si es válido)}
\ExtensionTok{python} \AttributeTok{{-}m}\NormalTok{ experiments.resolver }\AttributeTok{{-}{-}instance}\NormalTok{ data/test_instances/graph1.hcp }\AttributeTok{{-}{-}algorithm}\NormalTok{ aco\_ls }\AttributeTok{{-}{-}time{-}limit}\NormalTok{ 30}

\CommentTok{\# Benchmark principal y análisis (tablas, pruebas, métricas F1 y figuras)}
\ExtensionTok{python} \AttributeTok{{-}m}\NormalTok{ experiments.benchmark }\AttributeTok{{-}{-}instances}\NormalTok{ data/test_instances/graph1.hcp data/test_instances/graph24.hcp }\DataTypeTok{\textbackslash{}}
\NormalTok{    data/test_instances/graph41.hcp data/test_instances/graph172.hcp }\DataTypeTok{\textbackslash{}}
    \AttributeTok{{-}{-}algorithms}\NormalTok{ ga ga\_er sa sa\_2opt aco aco\_ls aco\_ls\_pre bee de pso random\_search }\DataTypeTok{\textbackslash{}}
    \AttributeTok{{-}{-}seeds}\NormalTok{ 10 }\AttributeTok{{-}{-}max{-}evals}\NormalTok{ 100000 }\AttributeTok{{-}{-}time{-}limit}\NormalTok{ 15 }\AttributeTok{{-}{-}out}\NormalTok{ results/bench }\AttributeTok{{-}{-}tours}\NormalTok{ results/tours}
\ExtensionTok{python} \AttributeTok{{-}m}\NormalTok{ experiments.tuning }\AttributeTok{{-}{-}out}\NormalTok{ results/tuning.csv}
\ExtensionTok{python} \AttributeTok{{-}m}\NormalTok{ experiments.tuning }\AttributeTok{{-}{-}out}\NormalTok{ results/tuning\_sa.csv }\AttributeTok{{-}{-}algs}\NormalTok{ sa }\AttributeTok{{-}{-}evals}\NormalTok{ 100000 }\DataTypeTok{\textbackslash{}}
    \AttributeTok{{-}{-}grid} \StringTok{\textquotesingle{}\{"t0": [0.1, 0.5, 1.0, 2.0, 3.0], "t\_end": [0.01, 0.02, 0.05, 0.1]\}\textquotesingle{}}
\ExtensionTok{python} \AttributeTok{{-}m}\NormalTok{ experiments.analisis }\AttributeTok{{-}{-}prefix}\NormalTok{ results/bench  }\AttributeTok{{-}{-}out}\NormalTok{ results/analisis}
\ExtensionTok{python} \AttributeTok{{-}m}\NormalTok{ experiments.analisis }\AttributeTok{{-}{-}prefix}\NormalTok{ results/pesada }\AttributeTok{{-}{-}out}\NormalTok{ results/analisis\_pesada}
\end{Highlighting}
\end{Shaded}

Los ciclos hamiltonianos encontrados para graph1 (20 en
\texttt{results/tours/} y 25 en \texttt{results/tours\_pesada/}) están
en formato TOUR del FHCP.

\begin{center}\rule{0.5\linewidth}{0.5pt}\end{center}

\clearpage
\section{9. Verificación de requisitos de la
práctica}\label{verificaciuxf3n-de-requisitos-de-la-pruxe1ctica}

{\def\LTcaptype{none} % do not increment counter
\begin{longtable}[]{@{}
  >{\raggedright\arraybackslash}p{(\linewidth - 2\tabcolsep) * \real{0.5000}}
  >{\raggedright\arraybackslash}p{(\linewidth - 2\tabcolsep) * \real{0.5000}}@{}}
\toprule\noalign{}
\begin{minipage}[b]{\linewidth}\raggedright
Requisito de \texttt{practica.pdf}
\end{minipage} & \begin{minipage}[b]{\linewidth}\raggedright
Dónde se cumple
\end{minipage} \\
\midrule\noalign{}
\endhead
\bottomrule\noalign{}
\endlastfoot
F1: Teoría de grafos y ciclos hamiltonianos (longitud n) & §1.1 \\
F1: Representación del espacio de búsqueda sin estructuras ilegales &
§1.2, §2.2 \\
F1: Decodificación, fitness y heurísticas & §1.3, §2.2 \\
F1: Métricas (costo, iteraciones, precisión, recall, F1) & §1.4, §5.2,
§5.4 \\
F1: Estándares, formatos y criterios del FHCP & §1.5 \\
F2: Problema del ciclo hamiltoniano, NP-difícil, relación con TSP &
§1.1 \\
F2: Dataset FHCP Challenge Set & §4.1 \\
F3: GA (base de control) & §3.1 \\
F3: SA con swaps y enfriamiento lento & §3.2 \\
F3: ACO y Bee Colony (cooperación de enjambre) & §3.3, §3.4 \\
F3: DE adaptada a permutaciones & §3.5 \\
F3: PSO & §3.6 \\
F4.1: Código que carga grafos FHCP y devuelve un arreglo de vértices &
\texttt{src/}, \texttt{experiments/resolver.py}, §2.3 \\
F4.2: Definición algorítmica del grafo y fundamentación teórica & §1.1,
§2.1 \\
F4.2: Mapeo de representación y función de fitness & §2.2 \\
F4.2: Comparativa (tasa de éxito y tiempo) en tablas y gráficos & §5 \\
F4.2: Benchmarking contra el FHCP Challenge & §6 \\
F4.2: Conclusiones sobre la metaheurística que mejor se adaptó & §7 \\
\end{longtable}
}

\begin{center}\rule{0.5\linewidth}{0.5pt}\end{center}


\clearpage
\section{10. Guía explicada de la implementación y del uso}
\subsection{Del archivo de entrada a una solución validada}
El recorrido del programa es: leer el archivo HCP, construir el grafo, seleccionar el algoritmo, generar candidatos, evaluarlos, conservar el mejor y validarlo. El lector está en \texttt{src/core/hcp\_io.py}; el modelo del grafo, en \texttt{src/core/graph.py}; y la evaluación y validación, en \texttt{src/core/fitness.py}. El registro de \texttt{src/algorithms/\_\_init\_\_.py} permite elegir la implementación por nombre.

La configuración común fija evaluaciones máximas, tiempo límite y parámetros específicos. El resultado contiene la mejor ruta, su costo, el historial de mejoras, el tiempo y las evaluaciones. El presupuesto se consulta durante la búsqueda; no es un corte externo exacto, por lo que una operación en curso puede producir un pequeño exceso de tiempo. La misma semilla hace reproducibles las decisiones aleatorias, pero un corte por tiempo puede cambiar la cantidad de trabajo ejecutado según el equipo.

\subsection{Ejemplo de representación, costo y cierre}
En el grafo de seis vértices \texttt{envelope6.hcp}, la ruta externa $(1,2,3,6,5,4)$ utiliza las aristas $1\!:\!2$, $2\!:\!3$, $3\!:\!6$, $6\!:\!5$, $5\!:\!4$ y $4\!:\!1$. Todas existen y el costo es cero. El arreglo contiene seis posiciones; no se agrega otro 1 al arreglo porque el regreso al origen es implícito. Internamente se almacena $(0,1,2,5,4,3)$.

La ruta $(1,2,3,4,5,6)$ también es una permutación legal, pero falta la arista $6\!:\!1$: su costo es uno. Esto distingue \emph{legalidad de la representación} de \emph{validez del ciclo}. Para $n=6$, la fracción de aristas existentes es $5/6$. Esa cercanía no basta para declarar éxito.

\subsection{Cómo trabaja cada método, paso a paso}
\textbf{GA.} Inicializa una población de permutaciones; selecciona padres por torneo; aplica cruza OX para conservar una parte del orden sin repetir vértices; muta por inversión; evalúa descendientes; y conserva élites. La población aporta diversidad y el elitismo protege las mejores rutas. Sin embargo, conservar orden no implica conservar aristas, por lo que una cruza puede deshacer conexiones útiles. La variante ER intenta heredar conexiones entre vértices.

\textbf{SA.} Comienza con una ruta y propone intercambiar dos posiciones. Si el costo disminuye o permanece igual, acepta el cambio; si aumenta en $\Delta$, lo acepta con probabilidad $\exp(-\Delta/T)$. A temperatura alta explora; al enfriarse se vuelve más selectivo. Por ejemplo, con $\Delta=1$ y $T=2$, la probabilidad es aproximadamente 0.607; con $T=0.01$, es prácticamente cero. Se conserva aparte el mejor candidato para no perderlo al aceptar un empeoramiento. La variante 2-opt invierte un segmento y cambia otro vecindario, sin sustituir al SA con swaps exigido.

\textbf{ACO.} Varias hormigas construyen rutas. Entre vecinos disponibles se combinan la feromona, que recoge experiencia compartida, y una preferencia heurística por vértices con pocas opciones libres. Si una hormiga queda bloqueada, completa mediante un salto penalizado. Tras construir rutas, la evaporación reduce memoria antigua y el depósito refuerza conexiones prometedoras. La variante con búsqueda local intenta corregir aristas faltantes antes de actualizar el mejor resultado.

\textbf{Colonia de abejas.} Cada fuente representa una ruta y su costo mide su calidad. Las empleadas modifican su fuente; las observadoras dedican más intentos a fuentes de mejor costo; las exploradoras reemplazan fuentes estancadas. La información compartida es la calidad de las fuentes. Esta implementación ABC no usa una matriz de feromonas: ese mecanismo corresponde a ACO. Su cooperación se basa en selección probabilística y abandono, una distinción importante al interpretar el enunciado.

\textbf{DE.} Representa cada candidato con un vector de claves reales. Combina tres vectores mediante una diferencia escalada y aplica cruza binomial; después ordena las claves para obtener la permutación y evalúa esa ruta. Si mejora o empata al padre, lo reemplaza. Por ejemplo, las claves $(0.3,0.1,0.4,0.2)$ se decodifican como los índices $(1,3,0,2)$. El ordenamiento garantiza que cada índice aparezca una sola vez, aunque puede transformar cambios pequeños de claves en cambios grandes de conexiones.

\textbf{PSO.} Cada partícula mantiene claves, velocidad, su mejor posición y la mejor conocida por el enjambre. La inercia conserva parte del movimiento; el componente cognitivo atrae al mejor personal; el social atrae al mejor compartido. Cada nueva posición se ordena para producir una permutación. El mejor global es una referencia compartida, no un agente que ejecute las decisiones de las demás partículas. La adaptación conserva rutas legales, pero no garantiza buena localidad en el espacio de aristas.

\subsection{Cómo interpretar las métricas sin confundirlas}
La tasa de éxito responde qué proporción de corridas encontró un ciclo completo; el fitness responde cuántas aristas faltaron. El tiempo mide el costo observado en el entorno de ejecución. Reportar las tres evita considerar ganador a un método que termina rápido pero no resuelve, o a uno que reduce el costo sin llegar a cero.

Las fórmulas de clasificación son $P=TP/(TP+FP)$, $R=TP/(TP+FN)$ y $F_1=2PR/(P+R)$. Si no se reporta ningún positivo, la precisión no está definida y se indica como no disponible. Si hay positivos reales y no se encuentra ninguno, el recall es cero; aquí se usa F1 igual a cero para ese caso. Al ser todas las instancias hamiltonianas, no se estima especificidad ni capacidad para reconocer grafos no hamiltonianos. Las pruebas estadísticas complementan los tamaños de efecto, sin convertir cuatro instancias en una muestra representativa de todo el FHCP.

Contar evaluaciones aproxima el esfuerzo, pero no lo iguala: una evaluación incremental de SA puede ser mucho más barata que construir y refinar una ruta en ACO. La búsqueda local realiza trabajo adicional dentro de una hormiga. Por ello se deben leer juntas las columnas de evaluaciones y tiempo. La ronda pesada no debe mezclarse con la principal como si usaran el mismo presupuesto.

\subsection{Instalación y primera ejecución}
Se requiere Python 3.10 o posterior por la sintaxis de tipos utilizada. Desde la raíz del proyecto:
\begin{Verbatim}[fontsize=\small,breaklines=true]
python3 -m venv .venv
source .venv/bin/activate
pip install -r requirements.txt
python -m experiments.interfaz
\end{Verbatim}
En Linux, si falta el módulo de interfaz, se instala con \texttt{sudo apt install python3-tk}. La interfaz necesita una sesión gráfica. Seleccionar \texttt{graph1.hcp}, elegir ACO + 2-opt, dejar semilla 0, 100000 evaluaciones y 15 segundos, y pulsar Ejecutar. La búsqueda se realiza en un hilo para mantener la ventana disponible. El panel muestra validación, costo, tiempo y arreglo; las gráficas muestran el ciclo propuesto y la convergencia.

El verde indica una arista existente; el rojo, una conexión faltante. La distribución circular es un dibujo de las etiquetas, no una representación de distancias geográficas. Una línea de convergencia horizontal significa que no mejoró el mejor costo durante ese intervalo; no implica que el algoritmo haya dejado de explorar. Guardar el TOUR únicamente cuando la validación confirma un ciclo completo.

Para ejecutar sin ventana y mostrar todos los vértices:
\begin{Verbatim}[fontsize=\small,breaklines=true]
python -m experiments.resolver \
  --instance data/test_instances/graph1.hcp \
  --algorithm aco_ls --seed 0 \
  --time-limit 30 --completo
\end{Verbatim}
El comando imprime etiquetas de 1 a $n$. Si el resultado es válido, guarda por defecto el TOUR en \texttt{results/tours/}; termina con código 0 si hay ciclo y con 1 si no lo encuentra. Este último caso es un resultado de búsqueda y no necesariamente un error del programa.

\subsection{Dataset completo, resultados y pruebas}
La carpeta \texttt{data/instances/} está vacía en la copia revisada; contiene únicamente un archivo de indicación. Para la demostración se usan los archivos incluidos en \texttt{data/test\_instances/}. Para experimentar con todo el set se debe extraer \texttt{data/FHCPCS.7z}, localizar los archivos HCP y colocarlos en \texttt{data/instances/}, o usar sus rutas reales. Tener un catálogo de 1001 grafos no significa haber ejecutado 1001 pruebas.

La procedencia local declara una copia pública del conjunto. La presencia de un archivo llamado FHCPCS.7z y la coincidencia de características no certifican por sí solas una descarga oficial. Debe conservarse la URL de descarga y comparar contenidos normalizando saltos de línea si se desea cerrar esa verificación. Este informe mantiene explícita esa limitación.

Los CSV de corridas están en \texttt{results/bench\_runs.csv} y \texttt{results/pesada\_runs.csv}; los JSONL conservan historiales; las tablas y figuras se encuentran en \texttt{results/analisis/} y \texttt{results/analisis\_pesada/}. Esos datos respaldan las comparativas; no se presentan las corridas históricas como experimentos repetidos durante esta edición.

La revisión ejecutó \texttt{python3 -m pytest -q}: 48 pruebas aprobaron y una se omitió porque busca instancias en la carpeta vacía. La prueba de uso de \texttt{resolver} con graph1 y ACO + 2-opt encontró un ciclo válido. La importación de la interfaz fue correcta; esto no sustituye una prueba visual completa de la ventana.

\subsection{Límites y mejoras del diseño experimental}
El resultado de una instancia resuelta entre seis probadas describe únicamente esa muestra. No permite afirmar una tasa de $1/1001$, ni extrapolar $1/6$ a todo el set. El resultado oficial de 985 grafos usa otra cobertura, otros recursos y técnicas especializadas; el contraste de calidad es contextual, no una competencia de velocidad bajo condiciones idénticas.

Para fortalecer una evaluación futura conviene seleccionar instancias por tamaño y estructura, fijar el hardware y ejecutar secuencialmente para comparar tiempos, separar instancias de ajuste y evaluación, y conservar versiones y semillas. Los presupuestos deben ser iguales dentro de cada experimento. Las conclusiones actuales se restringen a los datos observados: ACO con búsqueda local destaca en graph1, pero su sobrecosto puede perjudicarlo en grafos grandes.

\subsection{Correspondencia del portafolio con la práctica}
La sección 1 cubre la investigación; la 2, el grafo y el fitness; la 3, las implementaciones; la 4, el dataset y protocolo; la 5, las comparativas; la 6, el contraste oficial; y la 7, las conclusiones. Las secciones 8 y 9 permiten reproducir y localizar los requisitos. El portafolio incluye este informe, el código de \texttt{src/} y \texttt{experiments/}, y las evidencias de \texttt{results/}.

\clearpage
\begingroup\small\setstretch{1.0}
\section{Referencias}\label{referencias}

\begin{itemize}
\tightlist
\item
  Auger, A. y Hansen, N. (2005). Performance evaluation of an advanced
  local search evolutionary algorithm. \emph{IEEE CEC 2005}.
\item
  Bean, J. C. (1994). Genetic algorithms and random keys for sequencing
  and optimization. \emph{ORSA Journal on Computing}, 6(2).
\item
  Croes, G. A. (1958). A method for solving traveling-salesman problems.
  \emph{Operations Research}, 6(6).
\item
  Dorigo, M. y Stützle, T. (2004). \emph{Ant Colony Optimization}. MIT
  Press.
\item
  Garey, M. R. y Johnson, D. S. (1979). \emph{Computers and
  Intractability}. W. H. Freeman.
\item
  Goldberg, D. E. (1989). \emph{Genetic Algorithms in Search,
  Optimization and Machine Learning}. Addison-Wesley.
\item
  Haythorpe, M. (2018). FHCP Challenge Set: The first set of
  structurally difficult instances of the Hamiltonian cycle problem.
  \emph{Bulletin of the ICA}, 83. arXiv:1902.10352.
\item
  Helsgaun, K. (2000). An effective implementation of the Lin-Kernighan
  traveling salesman heuristic. \emph{European Journal of Operational
  Research}, 126(1).
\item
  Holm, S. (1979). A simple sequentially rejective multiple test
  procedure. \emph{Scandinavian Journal of Statistics}, 6(2).
\item
  Karaboga, D. y Basturk, B. (2007). A powerful and efficient algorithm
  for numerical function optimization: artificial bee colony (ABC)
  algorithm. \emph{Journal of Global Optimization}, 39(3).
\item
  Karp, R. M. (1972). Reducibility among combinatorial problems. En
  \emph{Complexity of Computer Computations}. Plenum.
\item
  Kennedy, J. y Eberhart, R. (1995). Particle swarm optimization.
  \emph{Proc. IEEE ICNN}.
\item
  Kirkpatrick, S., Gelatt, C. D. y Vecchi, M. P. (1983). Optimization by
  simulated annealing. \emph{Science}, 220(4598).
\item
  Pósa, L. (1976). Hamiltonian circuits in random graphs. \emph{Discrete
  Mathematics}, 14(4).
\item
  Storn, R. y Price, K. (1997). Differential evolution -- a simple and
  efficient heuristic for global optimization over continuous spaces.
  \emph{Journal of Global Optimization}, 11(4).
\item
  Stützle, T. y Hoos, H. H. (2000). MAX-MIN Ant System. \emph{Future
  Generation Computer Systems}, 16(8).
\item
  Tasgetiren, M. F., Liang, Y.-C., Sevkli, M. y Gencyilmaz, G. (2007). A
  particle swarm optimization algorithm for makespan and total flowtime
  minimization in the permutation flowshop sequencing problem.
  \emph{European Journal of Operational Research}, 177(3).
\item
  Whitley, D., Starkweather, T. y Fuquay, D. (1989). Scheduling problems
  and traveling salesmen: the genetic edge recombination operator.
  \emph{Proc. 3rd ICGA}.
\item
  Flinders Hamiltonian Cycle Project -- FHCP Challenge Set:
  \url{https://sites.flinders.edu.au/flinders-hamiltonian-cycle-project/fhcp-challenge-set/}
\end{itemize}


\paragraph{Fuentes verificadas para esta edición.}
Flinders University, \emph{FHCP Challenge Set}, página oficial, consultada el 13 de noviembre de 2026: \url{https://sites.flinders.edu.au/flinders-hamiltonian-cycle-project/fhcp-challenge-set/}. Se verificaron el tamaño del conjunto, los formatos HCP y TOUR, el criterio de competencia y los 985 grafos resueltos por los ganadores. Haythorpe, M., \emph{FHCP Challenge Set}, versión del artículo en \url{https://arxiv.org/abs/1902.10352}, publicación en Bulletin of the ICA, 83, pp. 98--107, 2018. La práctica original se conserva en \texttt{practica.pdf}; no es una fuente de resultados experimentales.

\endgroup
\end{document}
